\documentclass[11pt,a4paper]{article}

% ── Packages ──────────────────────────────────────────────────────────
\usepackage{amsmath,amssymb,amsthm}
\usepackage[margin=1in]{geometry}
\usepackage{hyperref}
\usepackage{graphicx}
\usepackage{booktabs}
\usepackage{xcolor}

% ── Theorem environments ──────────────────────────────────────────────
\newtheorem{theorem}{Theorem}
\newtheorem{lemma}[theorem]{Lemma}
\newtheorem{proposition}[theorem]{Proposition}
\newtheorem{corollary}[theorem]{Corollary}
\newtheorem{remark}[theorem]{Remark}
\newtheorem{definition}[theorem]{Definition}

% ── Operators ─────────────────────────────────────────────────────────
\DeclareMathOperator{\gap}{gap}
\DeclareMathOperator{\tr}{Tr}

% ── Metadata ──────────────────────────────────────────────────────────
\title{Lipschitz Stability of the Experiential Density Functional}
\author{Bryan Ehrlich}
\date{March 2026}

\begin{document}
\maketitle

% ======================================================================
% Abstract
% ======================================================================
\begin{abstract}
The experiential density $\rho(P) = I_\pi(B;M)\bigl(1 - I_\pi(B;M)/H_\pi(B)\bigr)$
measures meaningful complexity in finite-state Markov chains whose state
space decomposes as $\Omega = B \times M$.  We prove that $\rho$ is
Lipschitz continuous under kernel perturbations, with an explicit constant
$L = (C_I + C_H)/\gap(P)$ depending on the state-space dimension and
spectral gap.  The proof composes Cho--Meyer perturbation theory for
stationary distributions, Fannes--Audenaert entropy continuity, and the
multivariate mean value theorem applied to the product functional
$f(I,H) = I(1 - I/H)$.  We also state a non-linear (sub-additive) bound
that is sharper for small perturbations.  Numerical verification against
3000 random kernel perturbations of a 16-state observer chain confirms
$L_{\mathrm{numerical}} \le L_{\mathrm{proven}}$ with zero violations
and a tightness gap of $450$--$9500\times$, indicating the bound is
conservative but valid.  The proven bound's $1/\gap$ dependence is approximately consistent
with numerical observations ($R^2 = 0.97$, fitted exponent $-0.89$).
\end{abstract}

% ======================================================================
% 1. Introduction
% ======================================================================
\section{Introduction}

\subsection{The experiential density}

Consider a finite Markov chain on a composite state space
$\Omega = B \times M$, where $B$ represents the ``body'' (environment
or physical degrees of freedom) and $M$ represents the ``model''
(internal or representational degrees of freedom).  This decomposition
arises naturally in the experiential measure
framework~\cite{Ehrlich2026a}, which studies how information-processing
systems maintain internal models of their environment through dynamical
coupling.

Let $P$ be an irreducible row-stochastic matrix on $\Omega$ with
unique stationary distribution $\pi$.  Reshaping $\pi$ as a joint
distribution $p(b,m)$ on $B \times M$, the mutual information
$I_\pi(B;M) = H_\pi(B) + H_\pi(M) - H_\pi(B,M)$ quantifies the
statistical coupling between body and model in the stationary state,
and $H_\pi(B)$ is the marginal entropy of the body.

The \emph{experiential density} is defined as
\begin{equation}\label{eq:rho-intro}
  \rho(P) = I_\pi(B;M)\left(1 - \frac{I_\pi(B;M)}{H_\pi(B)}\right).
\end{equation}
This functional vanishes at two extremes: when $I = 0$ (body and model
are independent, no representation exists) and when $I = H(B)$ (the
model is a perfect copy of the body, leaving no room for further
experience).  It is maximized at $I = H(B)/2$, where the system balances
knowledge against uncertainty---a regime of maximal ``meaningful
complexity.''

The experiential density belongs to the LMC (Lopez-Ruiz--Mancini--Calbet)
statistical complexity family~\cite{LMC1995}, which measures complexity as
the product of information content and disequilibrium.  It is one member
of a broader family of complexity measures that share this multiplicative
structure; the Lipschitz stability techniques developed here (composing
stationary-distribution perturbation bounds with entropy continuity)
could be applied to derive analogous results for other members of the
LMC family.  More broadly, the experiential density instantiates the
effective complexity principle of Gell-Mann and
Lloyd~\cite{GellMannLloyd2003}: meaningful structure peaks between
pure order and pure randomness.

\subsection{Why stability matters}

A well-posed functional should exhibit continuous dependence on its
parameters.  If $\rho$ were discontinuous in the transition kernel
$P$---if an arbitrarily small change in dynamics could produce a finite
jump in experiential density---then $\rho$ would be poorly behaved as a
dynamical functional.  Lipschitz continuity is the strongest standard form of
continuity: it provides a linear bound on the output perturbation in
terms of the input perturbation, with an explicit constant.

Stability results are also prerequisite for any perturbation theory.
If we wish to study how $\rho$ responds to environmental noise, parameter
drift, or approximate computation, we need quantitative control over
$|\rho(P) - \rho(P')|$ in terms of $\|P - P'\|$.

\subsection{Contributions of this paper}

We prove that $\rho$ is Lipschitz continuous under kernel perturbations
(Theorem~\ref{thm:lipschitz}), with:
\begin{enumerate}
  \item A \textbf{non-linear bound} (sharpest form):
    $|\rho(P) - \rho(P')| \le
    \frac{\delta}{2}[2\ln(|B|-1) + \ln(|M|-1) + \ln(n-1)]
    + 4\,h_{\mathrm{bin}}(\delta/2)$,
    where $\delta = \|\pi - \pi'\|_1 \le \|P-P'\|_\infty/\gap(P)$.
  \item A \textbf{linear Lipschitz bound}:
    $|\rho(P) - \rho(P')| \le L(\delta)\cdot\|P-P'\|_\infty$
    with $L = (C_I + C_H)/\gap(P)$.
  \item \textbf{Explicit dependence} on state-space dimension
    ($\sim\ln|\Omega|$) and spectral gap ($\sim 1/\gap$), with
    no dependence on a lower bound for $H_\pi(B)$.
\end{enumerate}

The proof composes three standard results: Cho--Meyer perturbation
theory~\cite{ChoMeyer2001} (kernel $\to$ stationary distribution),
Fannes--Audenaert entropy continuity~\cite{Audenaert2007} (distribution
$\to$ information quantities), and the multivariate mean value theorem
(information quantities $\to$ $\rho$).  Each step yields an explicit
constant; their composition produces the final bound.

Using the mean value theorem rather than a direct product-rule
decomposition yields $L = (C_I + C_H)/\gap$ instead of the weaker
$(2C_I + C_H)/\gap$, because $|\partial f/\partial I| \le 1$ on the
feasible domain (not $\le 2$ as a naive product rule suggests).

Section~\ref{sec:numerical} reports numerical verification against
3000 random kernel perturbations.  The proven bound is never violated,
confirming correctness.  The tightness gap (ratio $L_{\mathrm{proven}}/
L_{\mathrm{numerical}}$) ranges from $450\times$ to $9500\times$,
indicating that while the bound is conservative, tighter results would
require fundamentally different techniques (e.g., direct differentiation
of $\rho$ as a function of $P$).


% ======================================================================
% 2. Setup and Statement
% ======================================================================
\section{Setup and Theorem Statement}\label{sec:setup}

\begin{definition}[Experiential density]\label{def:rho}
Let $P$ be an irreducible row-stochastic matrix on $\Omega = B \times M$
with stationary distribution $\pi$ (the unique row vector satisfying
$\pi P = \pi$, $\pi_i \ge 0$, $\sum_i \pi_i = 1$).
Reshape $\pi$ as a joint distribution $p(b,m)$ on $B \times M$.
The \emph{experiential density} is
\begin{equation}\label{eq:rho-def}
  \rho(P) = I_\pi(B;M)\left(1 - \frac{I_\pi(B;M)}{H_\pi(B)}\right),
\end{equation}
where $I_\pi(B;M) = H_\pi(B) + H_\pi(M) - H_\pi(B,M)$ is the mutual
information and $H_\pi(B) = -\sum_b p_B(b)\ln p_B(b)$ is the marginal
entropy on $B$, all in nats.
\end{definition}

\begin{definition}[Spectral gap and norms]\label{def:gap}
For an irreducible stochastic matrix $P$ on $n = |\Omega|$ states:
\begin{itemize}
  \item The \emph{absolute spectral gap} is
        $\gap(P) = 1 - |\lambda_2(P)|$, where $\lambda_2$ is the
        eigenvalue of $P$ with second-largest absolute value.
  \item The \emph{row-sum norm} is
        $\|P\|_\infty = \max_i \sum_j |P_{ij}|$.
        For the difference of stochastic matrices,
        $\|P - P'\|_\infty = \max_i \sum_j |P_{ij} - P'_{ij}|
        = \max_i \|P_i - P'_i\|_1$.
  \item We use the $\ell^1$ norm throughout:
        $\|p - q\|_1 = \sum_i |p_i - q_i|$.
        The total variation distance is
        $\|p-q\|_{\mathrm{TV}} = \tfrac{1}{2}\|p-q\|_1$.
\end{itemize}
\end{definition}

\begin{theorem}[Lipschitz stability of $\rho$]\label{thm:lipschitz}
Let $P$, $P'$ be irreducible row-stochastic matrices on
$\Omega = B \times M$ with $|B| \ge 2$, $|M| \ge 2$, $|\Omega| = n$.
Let $\pi$, $\pi'$ be their stationary distributions, and suppose
$\gap(P) > 0$.

\medskip\noindent
\textbf{Non-linear bound (sharpest form).}
Define $\delta = \|\pi - \pi'\|_1$.  Then:
\begin{equation}\label{eq:rho-nonlinear}
  |\rho(P) - \rho(P')| \le
  \frac{\delta}{2}\bigl[2\ln(|B|-1) + \ln(|M|-1) + \ln(n-1)\bigr]
  + 4\,h_{\mathrm{bin}}\!\left(\frac{\delta}{2}\right),
\end{equation}
where $h_{\mathrm{bin}}(x) = -x\ln x - (1-x)\ln(1-x)$ and
$\delta \le \|P - P'\|_\infty / \gap(P)$ by the Cho--Meyer bound.

\medskip\noindent
\textbf{Linear Lipschitz bound.}
Equivalently,
\begin{equation}\label{eq:lipschitz-bound}
  |\rho(P) - \rho(P')| \le L(\delta) \cdot \|P - P'\|_\infty,
\end{equation}
where
\begin{equation}\label{eq:L-explicit}
  L(\delta) = \frac{C_I(\delta) + C_H(\delta)}{\gap(P)},
\end{equation}
with
\begin{align}
  C_I(\delta) &= \frac{1}{2}\bigl[\ln(|B|-1) + \ln(|M|-1) + \ln(n-1)\bigr]
         + 3\,\frac{h_{\mathrm{bin}}(\delta/2)}{\delta}, \label{eq:CI} \\
  C_H(\delta) &= \frac{1}{2}\ln(|B|-1) + \frac{h_{\mathrm{bin}}(\delta/2)}{\delta}. \label{eq:CH}
\end{align}

For small perturbations ($\delta \ll 1$),
$h_{\mathrm{bin}}(\delta/2)/\delta \sim \frac{1}{2}\ln(2/\delta)$,
so the leading behavior is:
\begin{equation}\label{eq:L-leading}
  L \sim \frac{1}{2\,\gap(P)}\,
  \bigl[2\ln(|B|-1) + \ln(|M|-1) + \ln(n-1)\bigr]
  + \frac{2\ln(2/\delta)}{\gap(P)}.
\end{equation}

\medskip\noindent
\textbf{Uniform non-linear bound.}
Substituting $\delta \le \|P-P'\|_\infty/\gap(P)$ into~\eqref{eq:rho-nonlinear}
and using $h_{\mathrm{bin}}(x) \le \ln 2$ for all $x$:
\begin{equation}\label{eq:rho-uniform}
  |\rho(P) - \rho(P')| \le
  \frac{\|P-P'\|_\infty}{2\,\gap(P)}\bigl[2\ln(|B|-1) + \ln(|M|-1)
  + \ln(n-1)\bigr] + 4\ln 2.
\end{equation}

\medskip\noindent
\textbf{Domain of validity:}
\begin{enumerate}
  \item $P$ irreducible with $\gap(P) > 0$;
  \item $P'$ irreducible (so $\pi'$ exists and is unique);
  \item $|B| \ge 2$ (otherwise $H_\pi(B) = 0$ and $\rho \equiv 0$).
\end{enumerate}
No lower bound on $H_\pi(B)$ is required in the Lipschitz constant.
\end{theorem}

The proof proceeds in three steps, each invoking a standard result.


% ======================================================================
% 3. Step 1 -- Stationary Distribution Perturbation
% ======================================================================
\section{Step 1: Stationary Distribution Perturbation}\label{sec:step1}

\begin{lemma}[Cho--Meyer perturbation bound~\cite{ChoMeyer2001}]
\label{lem:cho-meyer}
Let $P$, $P'$ be irreducible stochastic matrices on $n$ states with
stationary distributions $\pi$, $\pi'$.  Then
\begin{equation}\label{eq:cho-meyer}
  \|\pi - \pi'\|_1 \le
  \frac{1}{\gap(P)}\,\|P - P'\|_\infty.
\end{equation}
\end{lemma}

\begin{proof}[Proof sketch (following Cho--Meyer~\cite{ChoMeyer2001})]
Let $\Pi = \mathbf{1}\pi$ be the rank-one stationary projector.
From $\pi P = \pi$ and $\pi' P' = \pi'$:
\[
  \pi' - \pi = \pi'(P' - P)(I - P + \Pi)^{-1},
\]
where $(I - P + \Pi)^{-1}$ is the group inverse of $(I - P)$.
The key estimate is $\|(I - P + \Pi)^{-1}\|_\infty \le 1/\gap(P)$
(Meyer~\cite{Meyer1975}; Cho--Meyer~\cite{ChoMeyer2001}, Theorem~3.1).
Since $\|\pi'\|_1 = 1$:
\[
  \|\pi' - \pi\|_1
  \le \|\pi'\|_1 \cdot \|P' - P\|_\infty \cdot \|(I - P + \Pi)^{-1}\|_\infty
  \le \frac{\|P - P'\|_\infty}{\gap(P)}.
\]
Here we used $\|v A\|_1 \le \|v\|_1 \cdot \|A\|_\infty$ for a row
vector $v$ and matrix $A$ (standard submultiplicativity for the
$\ell^1$-$\ell^\infty$ pair).
\end{proof}

\begin{remark}\label{rem:cho-meyer-tight}
The bound~\eqref{eq:cho-meyer} may be loose for specific chains.
The sharper Cho--Meyer condition number
$\kappa(P) = \|\pi\|_\infty \cdot \|(I - P + \Pi)^{-1}\|_\infty$
can be much smaller than $1/\gap(P)$.  We use the simpler bound because
it yields clean dependence on the spectral gap.
\end{remark}

\textbf{Interface to Step~2.}
Define $\delta = \|\pi - \pi'\|_1$.  Reshaping $\pi$ and $\pi'$ as
joint distributions $p$ and $q$ on $B \times M$ (with
$p(b,m) = \pi_{(b,m)}$), we have $\|p - q\|_1 = \delta$.
By Lemma~\ref{lem:cho-meyer}:
\begin{equation}\label{eq:delta-bound}
  \delta \le \frac{\|P - P'\|_\infty}{\gap(P)}.
\end{equation}


% ======================================================================
% 4. Step 2 -- Entropy and Mutual Information Continuity
% ======================================================================
\section{Step 2: Mutual Information and Entropy Continuity}\label{sec:step2}

\subsection{Entropy continuity lemma}

\begin{lemma}[Fannes--Audenaert inequality~\cite{CsiszarKorner2011,Audenaert2007}]
\label{lem:entropy-cont}
Let $p$, $q$ be probability distributions on an alphabet of size $k \ge 2$
with $\|p - q\|_1 = \delta \in (0, 2(1 - 1/k)]$.  Then
\begin{equation}\label{eq:fannes}
  |H(p) - H(q)| \le \frac{\delta}{2}\,\ln(k-1)
  + h_{\mathrm{bin}}\!\left(\frac{\delta}{2}\right),
\end{equation}
where $h_{\mathrm{bin}}(x) = -x\ln x - (1-x)\ln(1-x)$ in nats.
\end{lemma}

\begin{remark}[Convention]\label{rem:convention}
Many references state this bound with $\delta_{\mathrm{TV}} = \|p-q\|_{\mathrm{TV}}
= \frac{1}{2}\|p-q\|_1$.  In that convention, the bound reads
$|H(p) - H(q)| \le \delta_{\mathrm{TV}}\ln(k-1) + h_{\mathrm{bin}}(\delta_{\mathrm{TV}})$.
Our formula~\eqref{eq:fannes} is identical after substituting
$\delta_{\mathrm{TV}} = \delta/2$.
The tight constant is due to Audenaert~\cite{Audenaert2007}.
\end{remark}

\subsection{Marginal contraction}

\begin{lemma}[Marginal contraction]\label{lem:marginal}
Let $p$, $q$ be joint distributions on $B \times M$.  Then
\begin{equation}\label{eq:marginal-contract}
  \|p_B - q_B\|_1 \le \|p - q\|_1,
  \qquad
  \|p_M - q_M\|_1 \le \|p - q\|_1,
\end{equation}
where $p_B(b) = \sum_m p(b,m)$ and $p_M(m) = \sum_b p(b,m)$.
\end{lemma}

\begin{proof}
By the triangle inequality:
$\|p_B - q_B\|_1 = \sum_b |p_B(b) - q_B(b)|
= \sum_b \bigl|\sum_m (p(b,m) - q(b,m))\bigr|
\le \sum_b \sum_m |p(b,m) - q(b,m)| = \|p - q\|_1$.
\end{proof}

\subsection{Mutual information continuity bound}

Using $I(B;M) = H(B) + H(M) - H(B,M)$ and the triangle inequality:
\begin{equation}\label{eq:MI-triangle}
  |I_p(B;M) - I_q(B;M)| \le
  |H_p(B) - H_q(B)| + |H_p(M) - H_q(M)| + |H_p(B,M) - H_q(B,M)|.
\end{equation}

Apply Lemma~\ref{lem:entropy-cont} to each term.  The marginal
contraction (Lemma~\ref{lem:marginal}) gives
$\|p_B - q_B\|_1 \le \delta$ and $\|p_M - q_M\|_1 \le \delta$,
where $\delta = \|p - q\|_1$.

\begin{itemize}
  \item $H(B)$: alphabet size $|B|$, perturbation $\le \delta$:
    $|H_p(B) - H_q(B)| \le \frac{\delta}{2}\ln(|B|-1) + h_{\mathrm{bin}}(\delta/2)$.

  \item $H(M)$: alphabet size $|M|$, perturbation $\le \delta$:
    $|H_p(M) - H_q(M)| \le \frac{\delta}{2}\ln(|M|-1) + h_{\mathrm{bin}}(\delta/2)$.

  \item $H(B,M)$: alphabet size $n = |B| \cdot |M|$, perturbation $= \delta$:
    $|H_p(B,M) - H_q(B,M)| \le \frac{\delta}{2}\ln(n-1) + h_{\mathrm{bin}}(\delta/2)$.
\end{itemize}

\noindent Summing:
\begin{equation}\label{eq:MI-bound}
  |I_p(B;M) - I_q(B;M)| \le
  \frac{\delta}{2}\bigl[\ln(|B|-1) + \ln(|M|-1) + \ln(n-1)\bigr]
  + 3\,h_{\mathrm{bin}}\!\left(\frac{\delta}{2}\right).
\end{equation}

\subsection{Separate bound on $H(B)$}

For the composition in Step~3, we also need:
\begin{equation}\label{eq:HB-bound}
  |H_p(B) - H_q(B)| \le
  \frac{\delta}{2}\ln(|B|-1) + h_{\mathrm{bin}}\!\left(\frac{\delta}{2}\right).
\end{equation}

\begin{remark}
An alternative approach to bounding $|I_p - I_q|$ would use conditional
entropy continuity results, such as the tight bounds of
Winter~\cite{Winter2016}.  Since mutual information equals $H(B) - H(B|M)$,
a direct continuity bound on the conditional entropy could potentially
yield tighter constants by avoiding the three-term triangle inequality
decomposition~\eqref{eq:MI-triangle}.  We leave this direction as future
work.
\end{remark}


% ======================================================================
% 5. Step 3 -- Composition
% ======================================================================
\section{Step 3: Composition via the Mean Value Theorem}\label{sec:step3}

Write $\rho = f(I, H)$ where $f(I, H) = I(1 - I/H)$, with
$I = I_\pi(B;M) \ge 0$, $H = H_\pi(B) > 0$, and $0 \le I \le H$.

\subsection{Gradient bound}

The partial derivatives of $f$ on the domain
$\mathcal{D} = \{(I,H) : 0 \le I \le H,\, H > 0\}$ are:
\begin{align}
  \frac{\partial f}{\partial I} &= 1 - \frac{2I}{H},
  \qquad\text{so } \left|\frac{\partial f}{\partial I}\right| \le 1
  \quad\text{(since $I/H \in [0,1]$)}; \label{eq:df-dI}\\[4pt]
  \frac{\partial f}{\partial H} &= \frac{I^2}{H^2},
  \qquad\text{so } \left|\frac{\partial f}{\partial H}\right| \le 1
  \quad\text{(since $I/H \in [0,1]$)}. \label{eq:df-dH}
\end{align}

The domain $\mathcal{D}$ is convex: for $(I,H), (I',H') \in \mathcal{D}$
and $t \in [0,1]$, the convex combination satisfies
$\bar{H} = (1-t)H + tH' > 0$ and
$\bar{I} = (1-t)I + tI' \le (1-t)H + tH' = \bar{H}$, so
$(\bar{I},\bar{H}) \in \mathcal{D}$.

By the multivariate mean value theorem, there exists
$(\bar{I},\bar{H})$ on the line segment from $(I,H)$ to $(I',H')$
such that:
\begin{equation}\label{eq:mvt-exact}
  f(I,H) - f(I',H') = \frac{\partial f}{\partial I}\bigg|_{(\bar{I},\bar{H})}
  (I - I') + \frac{\partial f}{\partial H}\bigg|_{(\bar{I},\bar{H})}(H - H').
\end{equation}

Taking absolute values and applying the gradient bounds:
\begin{equation}\label{eq:mvt}
  |f(I,H) - f(I',H')|
  \le \left|\frac{\partial f}{\partial I}\right|\,|I - I'|
  + \left|\frac{\partial f}{\partial H}\right|\,|H - H'|
  \le |I - I'| + |H - H'|.
\end{equation}

\begin{remark}
This bound is clean: no factors of $H$, $H'$, or $h_{\min}$ appear.
The key observation is that $f(I,H) = I(1-I/H)$ satisfies
$\|\nabla f\|_\infty \le 1$ everywhere on its feasible domain,
because $I/H \in [0,1]$.  A naive product-rule decomposition would
give $|f - f'| \le |I - I'| \cdot (1 - I/H) + I' \cdot |I/H - I'/H'|$,
which leads to the weaker constant $2C_I + C_H$ instead of $C_I + C_H$.
\end{remark}

\subsection{Final composition}

Combining~\eqref{eq:mvt} with the mutual information
bound~\eqref{eq:MI-bound}, the entropy bound~\eqref{eq:HB-bound},
and the Cho--Meyer bound~\eqref{eq:delta-bound}:

\begin{align}
  |\rho(P) - \rho(P')|
  &\le |I - I'| + |H - H'| \notag\\
  &\le \left[\frac{\delta}{2}\bigl(\ln(|B|-1) + \ln(|M|-1) + \ln(n-1)\bigr)
       + 3\,h_{\mathrm{bin}}\!\left(\frac{\delta}{2}\right)\right] \notag\\
  &\quad + \left[\frac{\delta}{2}\ln(|B|-1)
       + h_{\mathrm{bin}}\!\left(\frac{\delta}{2}\right)\right] \notag\\
  &= \frac{\delta}{2}\bigl[2\ln(|B|-1) + \ln(|M|-1) + \ln(n-1)\bigr]
     + 4\,h_{\mathrm{bin}}\!\left(\frac{\delta}{2}\right),
     \label{eq:rho-bound-nonlinear}
\end{align}
with $\delta \le \|P - P'\|_\infty / \gap(P)$.
This completes the proof of the non-linear bound~\eqref{eq:rho-nonlinear}.

\medskip
\textbf{Linear Lipschitz form.}  Writing
$h_{\mathrm{bin}}(\delta/2) \le \delta \cdot [h_{\mathrm{bin}}(\delta/2)/\delta]$
and defining $\eta(\delta) = h_{\mathrm{bin}}(\delta/2)/\delta$ (finite for
any $\delta > 0$):
\begin{align}
  |\rho(P) - \rho(P')|
  &\le \bigl[C_I(\delta) + C_H(\delta)\bigr]\,\delta \notag\\
  &\le \frac{C_I(\delta) + C_H(\delta)}{\gap(P)}\,\|P - P'\|_\infty,
  \label{eq:rho-bound-linear}
\end{align}
where $C_I(\delta)$ and $C_H(\delta)$ are as in~\eqref{eq:CI}--\eqref{eq:CH}.
This gives the Lipschitz constant:
\begin{equation}\label{eq:L-final}
  \boxed{L(\delta) = \frac{C_I(\delta) + C_H(\delta)}{\gap(P)}\,.}
\end{equation}

\begin{remark}[On the $\delta$-dependence of $L$]\label{rem:delta-dep}
The Lipschitz ``constant'' $L(\delta)$ depends on $\delta$ through
the $h_{\mathrm{bin}}/\delta$ ratio, which grows as $\frac{1}{2}\ln(2/\delta)$
for small $\delta$.  This is not a deficiency of the proof: it reflects
a genuine feature of entropy continuity.  Shannon entropy is \emph{not}
globally Lipschitz on the probability simplex, because
$|H(p) - H(q)|/\|p-q\|_1$ diverges logarithmically as $p$ approaches
a vertex of the simplex.

For any \emph{fixed} perturbation size $\varepsilon = \|P - P'\|_\infty$,
the constant $L(\varepsilon/\gap(P))$ is finite and explicit.
For the non-linear bound~\eqref{eq:rho-nonlinear}, no $\delta$-dependent
constant is needed: it is valid for all $\delta$ simultaneously.
\end{remark}


% ======================================================================
% 6. Explicit Formulas
% ======================================================================
\section{Explicit Formulas}\label{sec:explicit}

For reference, we collect the main formulas.  Write $n = |B|\cdot|M|$,
$g = \gap(P)$, $\varepsilon = \|P - P'\|_\infty$,
$\delta = \|\pi - \pi'\|_1 \le \varepsilon/g$.

\medskip\noindent
\textbf{Non-linear bound (recommended for evaluation):}
\begin{equation}\label{eq:nonlinear-final}
  |\rho(P) - \rho(P')| \le
  \frac{\varepsilon}{2\,g}\bigl[2\ln(|B|-1) + \ln(|M|-1) + \ln(n-1)\bigr]
  + 4\,h_{\mathrm{bin}}\!\left(\frac{\varepsilon}{2\,g}\right).
\end{equation}

\noindent
\textbf{Linear bound (with $\delta$-dependent constant):}
\begin{equation}\label{eq:L-collected}
  L(\delta) = \frac{1}{g}\left\{
    \frac{1}{2}\bigl[2\ln(|B|-1) + \ln(|M|-1) + \ln(n-1)\bigr]
    + 4\,\frac{h_{\mathrm{bin}}(\delta/2)}{\delta}
  \right\}.
\end{equation}

\noindent
\textbf{Asymptotic form ($\delta \to 0$):}
\begin{equation}\label{eq:L-asymptotic}
  L \sim \frac{1}{2\,g}\bigl[2\ln(|B|-1) + \ln(|M|-1) + \ln(n-1)\bigr]
  + \frac{2}{g}\ln\!\left(\frac{2\,g}{\varepsilon}\right).
\end{equation}

\noindent
\textbf{Uniform additive bound (all perturbation sizes):}
\begin{equation}\label{eq:uniform}
  |\rho(P) - \rho(P')| \le
  \frac{\varepsilon}{2\,g}\bigl[2\ln(|B|-1) + \ln(|M|-1) + \ln(n-1)\bigr]
  + 4\ln 2,
\end{equation}
using $h_{\mathrm{bin}}(x) \le \ln 2$ for all $x \in [0,1]$.


% ======================================================================
% 7. Limiting Cases
% ======================================================================
\section{Limiting Cases}\label{sec:limits}

We verify the bound's behavior in six regimes.

\subsection{Zero perturbation ($P' = P$)}
When $P' = P$: $\varepsilon = 0$, $\delta = 0$,
and~\eqref{eq:rho-nonlinear} gives $|\rho - \rho'| \le 0$, which is correct.

\subsection{Fast mixing ($\gap(P) \to 1$)}
For $\gap(P) = 1$ (e.g., a rank-one matrix $P = \mathbf{1}\pi$):
\[
  L_{\mathrm{eff}} \sim \frac{1}{2}\bigl[2\ln(|B|-1) + \ln(|M|-1) + \ln(n-1)\bigr].
\]
For the toy model ($|B| = |M| = 4$, $n = 16$):
$L_{\mathrm{eff}} \approx 3.00$ nats.  Since $\rho_{\max} = H(B)/4
\le \ln(4)/4 \approx 0.347$ nats, a Lipschitz constant of ${\sim}3$
means the bound allows $\rho$ to change by its full range when
$\varepsilon \sim 0.1$, which is reasonable for the fast-mixing limit.

\subsection{Slow mixing ($\gap(P) \to 0$)}
As $\gap(P) \to 0$:
the Cho--Meyer bound gives
$\delta \le \varepsilon/\gap(P) \to \infty$ (capped at $\delta = 2$);
$L \to \infty$.
This is physically correct: nearly reducible chains have stationary
distributions that are extremely sensitive to small kernel perturbations,
as tiny changes can redirect stationary mass between barely-communicating
basins.

\subsection{Independent $B$, $M$ ($I(B;M) = 0$)}
When the stationary distribution factorizes ($p(b,m) = p_B(b)\,p_M(m)$),
$I(B;M) = 0$ and $\rho = 0$.  Near this boundary, $\rho \sim I$
(to leading order, since $1 - I/H(B) \to 1$).
The bound gives $|\rho - \rho'| \le |I - I'| + |H - H'|$, consistent
with small perturbations producing small $\rho'$.

\subsection{Perfect tracking ($I(B;M) = H(B)$)}
When $M$ is a deterministic function of $B$, $I(B;M) = H(B)$ and
$\rho = 0$.  Near this boundary, $\rho \sim H(B) - I = H(M|B)$.
The bound handles this via $|\partial_I f| = |1 - 2I/H| \to 1$ as
$I \to H$, so perturbations away from perfect tracking increase
$\rho$ linearly.

\subsection{Toy model evaluation}
For the 7-chain toy model ($|B| = |M| = 4$, $n = 16$), the logarithmic
coefficient is:
\[
  \frac{1}{2}\bigl[2\ln 3 + \ln 3 + \ln 15\bigr]
  = \frac{1}{2}[2.197 + 1.099 + 2.708] = 3.002.
\]
Representative bound values:
\begin{center}
\begin{tabular}{@{}cccl@{}}
\toprule
$\gap$ & $\varepsilon$ & $\delta$ & Non-linear bound (nats) \\
\midrule
1.0 & 0.01 & 0.01 & 0.156 \\
1.0 & 0.10 & 0.10 & 1.094 \\
0.5 & 0.01 & 0.02 & 0.284 \\
0.5 & 0.10 & 0.20 & 1.901 \\
0.1 & 0.01 & 0.10 & 1.094 \\
0.1 & 0.10 & 1.00 & 5.775 \\
\bottomrule
\end{tabular}
\end{center}
Since $\rho_{\max} \approx 0.347$ nats, the bound is non-vacuous
for small perturbations with moderate spectral gap.


% ======================================================================
% 8. Numerical Verification
% ======================================================================
\section{Numerical Verification}\label{sec:numerical}

We verify Theorem~\ref{thm:lipschitz} numerically against random
kernel perturbations.  All computations use the verification code
\texttt{code/lipschitz\_verification.py} with seed 20260316 for
reproducibility.

\subsection{Observer chain}

The test chain is a 16-state Markov chain on $\Omega = B \times M$
with $|B| = |M| = 4$.  The body $B$ follows a metastable two-basin
dynamics (basins $\{0,1\}$ and $\{2,3\}$ with rare cross-basin
transitions), and the model $M$ tracks the body state with accuracy
$\alpha = 0.5$ and noise floor $\gamma = 0.02$.  This is the ``7-chain
observer'' from the experiential measure framework.

The chain's stationary properties are:
\begin{center}
\begin{tabular}{@{}ll@{}}
\toprule
Parameter & Value \\
\midrule
$|B|$, $|M|$, $|\Omega|$ & 4, 4, 16 \\
$\gap(P)$ & 0.100 \\
$\rho(P)$ & 0.346 nats \\
$I_\pi(B;M)$ & 0.704 nats \\
$H_\pi(B)$ & 1.386 nats \\
$I/H$ & 0.508 \\
\bottomrule
\end{tabular}
\end{center}

The ratio $I/H \approx 0.5$ places the chain near the maximum of
$f(I,H) = I(1-I/H)$, which peaks at $I = H/2$.

\subsection{Protocol}

For each perturbation size $\varepsilon \in \{0.001, 0.01, 0.1\}$,
we generate 1000 random kernel perturbations $P'$ of $P$ with
$\|P - P'\|_\infty \approx \varepsilon$.  Each perturbation is
constructed by adding a zero-sum random vector to each row, clipping
to non-negative entries, and renormalizing.  Only irreducible
perturbations (with $\gap(P') > 10^{-6}$) are retained.

For each perturbation we compute:
\begin{itemize}
  \item The numerical Lipschitz ratio
    $R_i = |\rho(P) - \rho(P'_i)| / \|P - P'_i\|_\infty$.
  \item The proven bound
    $L_{\mathrm{proven}} =
    \text{(nonlinear bound at $\delta = \varepsilon/\gap$)} / \varepsilon$.
\end{itemize}

The test passes if $\max_i R_i \le L_{\mathrm{proven}}$ for all
$\varepsilon$ values.

\subsection{Results}

\begin{center}
\begin{tabular}{@{}clllr@{}}
\toprule
$\varepsilon$ & $L_{\mathrm{numerical}}$ & $L_{\mathrm{proven}}$ &
  Tightness ratio & Violations \\
\midrule
0.001 & 0.012 & 109.4 & $9500\times$ & 0 / 1000 \\
0.01  & 0.019 & 109.4 & $5900\times$ & 0 / 1000 \\
0.1   & 0.115 &  51.4 & $450\times$  & 0 / 1000 \\
\bottomrule
\end{tabular}
\end{center}

Across all 3000 perturbations, $L_{\mathrm{numerical}} \le
L_{\mathrm{proven}}$ with zero violations.  The bound is confirmed
valid.

Figure~\ref{fig:verification} shows the distribution of Lipschitz
ratios at $\varepsilon = 0.01$.  The numerical ratios cluster near
zero (median $\approx 0.005$), while $L_{\mathrm{proven}} = 109.4$ lies
far to the right.  This visual confirms both validity and conservatism.

\begin{figure}[t]
  \centering
  \includegraphics[width=\textwidth]{../../figures/lipschitz_verification.pdf}
  \caption{Four-panel verification of the Lipschitz bound.
  (a)~Distribution of $|\rho(P) - \rho(P')|/\|P-P'\|_\infty$ for 1000
    random perturbations at $\varepsilon = 0.01$; the proven bound
    $L_{\mathrm{proven}} = 109.4$ (dashed red) is never exceeded.
  (b)~Log-log plot of $L_{\mathrm{numerical}}$ vs.\ $\gap(P)$,
    showing approximate $L \sim \gap^{-0.89}$ scaling ($R^2 = 0.97$).
  (c)~Gap-normalized Lipschitz constant $L \cdot \gap$ vs.\
    $\ln|\Omega|$; the linear fit ($R^2 = 0.97$) is a consistency
    check of the proven formula's logarithmic structure (see text).
  (d)~Convergence of $L_{\mathrm{numerical}}$ with sample size,
    stabilizing by $N \approx 500$.}
  \label{fig:verification}
\end{figure}

\subsection{Scaling laws}

\paragraph{Spectral gap dependence.}
Varying the inter-basin transition rate $q_{\mathrm{cross}}$ produces
chains with spectral gaps ranging from 0.005 to 0.3.  A log-log
regression of $L_{\mathrm{numerical}}$ against $\gap(P)$ yields:
\[
  L_{\mathrm{numerical}} \sim c \cdot \gap(P)^{-0.89},
  \qquad R^2 = 0.97.
\]
The exponent $-0.89$ is approximately consistent with the theoretical
prediction of $-1$.  The deviation may reflect finite-sample
underestimation (since $L_{\mathrm{numerical}}$ measures the empirical
maximum over a finite sample rather than the true supremum),
Cho--Meyer looseness for random (non-adversarial) perturbations, or both.

\paragraph{State-space size dependence.}
Varying $|B| = |M| = k$ for $k \in \{2,3,4,5,6\}$ and plotting the
gap-normalized quantity $L_{\mathrm{proven}} \cdot \gap$ against
$\ln|\Omega|$ yields a linear fit with $R^2 = 0.97$.  We note, however,
that this test evaluates the proven formula against itself: $L_{\mathrm{proven}}$
contains logarithmic terms by construction, so this is a consistency
check of the formula's internal structure rather than an independent
confirmation of $\ln|\Omega|$ scaling in the empirical data.

\subsection{Convergence}

Drawing 2000 perturbations cumulatively and tracking the running maximum,
$L_{\mathrm{numerical}}$ stabilizes by $N \approx 500$ samples
(relative change $< 5\%$ between successive thresholds), confirming
that 1000 samples per $\varepsilon$ are sufficient for reliable
estimates.


% ======================================================================
% 9. Discussion
% ======================================================================
\section{Discussion}\label{sec:discussion}

\subsection{Tightness of the bound}

The tightness gap of $450$--$9500\times$ between
$L_{\mathrm{proven}}$ and $L_{\mathrm{numerical}}$ deserves comment.
Four sources of slack compound:

\begin{enumerate}
  \item[(a)] \textbf{Cho--Meyer step.}  The bound
    $\delta \le \varepsilon/\gap$ is tight for adversarial perturbations
    that push all stationary mass in the direction of maximum
    sensitivity.  Random perturbations average over many directions and
    typically produce $\delta \ll \varepsilon/\gap$.

  \item[(b)] \textbf{Entropy continuity.}  The Fannes--Audenaert bound
    is tight for worst-case distributions (supported on two points),
    not for the specific marginals of an observer chain with rich
    internal structure.

  \item[(c)] \textbf{Gradient bound.}  The MVT uses
    $|\partial_I f| \le 1$ and $|\partial_H f| \le 1$ (worst case
    over $I/H \in [0,1]$).  For the specific $I/H \approx 0.508$ of
    the test chain, $|\partial_I f| = |1 - 2(0.508)| = 0.016$ and
    $|\partial_H f| = 0.508^2 = 0.258$, giving an effective gradient
    much smaller than the worst case.

  \item[(d)] \textbf{Correlation structure.}  The decomposition
    $|\Delta\rho| \le |\Delta I| + |\Delta H|$ overcounts because $I$
    and $H$ are correlated (both depend on $\pi$), so their
    perturbations may partially cancel rather than add.
\end{enumerate}

\subsection{Towards tighter bounds}

A fundamentally tighter bound would bypass the three-step decomposition
and instead compute $d\rho/dP_{ij}$ directly.  Writing
$\rho(P) = f(I(\pi(P)), H(\pi(P)))$, the chain rule gives
\[
  \frac{\partial \rho}{\partial P_{ij}}
  = \frac{\partial f}{\partial I}\,\frac{\partial I}{\partial P_{ij}}
  + \frac{\partial f}{\partial H}\,\frac{\partial H}{\partial P_{ij}},
\]
where $\partial I/\partial P_{ij}$ and $\partial H/\partial P_{ij}$
can be expressed through $\partial \pi/\partial P_{ij}$, which involves
the group inverse $(I - P + \Pi)^{-1}$.  This approach would give the
exact Lipschitz constant for a specific chain, but at the cost of
chain-specific computation rather than a universal formula.

\subsection{Role in the experiential measure framework}

This stability result is the second of three foundational properties
needed for the experiential density to serve as a well-behaved
dynamical functional:

\begin{enumerate}
  \item \textbf{Existence and uniqueness} (Theorem~A): $\rho(P)$ is
    well-defined for any irreducible $P$ on $B \times M$.
  \item \textbf{Lipschitz stability} (this paper): $\rho$ depends
    continuously on the kernel, with quantitative control.
  \item \textbf{Measurability}: $\rho(P)$ can in principle be estimated
    from finite trajectory data.
\end{enumerate}

The Lipschitz bound also provides error propagation: if the kernel $P$
is known only to accuracy $\varepsilon$ (e.g., from finite-data
estimation), then $\rho(P)$ is known to accuracy at most
$L(\varepsilon/\gap) \cdot \varepsilon$.  For the 7-chain observer
with $\gap = 0.1$ and $\varepsilon = 0.01$, the non-linear
bound~\eqref{eq:nonlinear-final} gives $|\Delta\rho| \le 1.09$ nats,
which exceeds $\rho_{\max} \approx 0.347$ and is therefore vacuous.
In practice, the actual perturbation is much smaller
($|\Delta\rho| \sim 10^{-4}$ nats), confirming that the bound is
conservative but that $\rho$ is genuinely stable.

More generally, the bound~\eqref{eq:nonlinear-final} is non-vacuous
(i.e., smaller than $\rho_{\max} = (\ln|B|)^2/(4\ln|B|) = \ln|B|/4$)
when the perturbation $\varepsilon$ is small enough relative to the
spectral gap and state-space size.  As a rough criterion, the leading
term dominates for small $\varepsilon$, so the bound is non-vacuous when
\[
  \frac{\varepsilon}{2\,\gap}\bigl[2\ln(|B|-1) + \ln(|M|-1) + \ln(n-1)\bigr]
  \lesssim \frac{\ln|B|}{4}\,,
\]
i.e., when $\varepsilon \lesssim \gap \cdot \ln|B| / (8\ln n)$.  For
chains with moderate spectral gap ($\gap \sim 0.1$) and modest state
space ($n \sim 16$), this requires $\varepsilon \lesssim 0.006$,
confirming that the bound is informative only for quite small
perturbations of such chains.


% ======================================================================
% Bibliography
% ======================================================================
\begin{thebibliography}{99}

\bibitem{LMC1995}
R.~Lopez-Ruiz, H.~L.~Mancini, X.~Calbet,
``A statistical measure of complexity,''
\emph{Phys.\ Lett.\ A} \textbf{209} (1995) 321--326.

\bibitem{GellMannLloyd2003}
M.~Gell-Mann, S.~Lloyd,
``Effective complexity,''
in \emph{Nonextensive Entropy}, eds.\ M.~Gell-Mann, C.~Tsallis,
Oxford Univ.\ Press (2004); original preprint 1996/2003.

\bibitem{ChoMeyer2001}
G.~E.~Cho, C.~D.~Meyer,
``Comparison of perturbation bounds for the stationary distribution
of a Markov chain,''
\emph{Linear Algebra Appl.} \textbf{335} (2001) 137--150.

\bibitem{Meyer1975}
C.~D.~Meyer,
``The role of the group generalized inverse in the theory of
finite Markov chains,''
\emph{SIAM Rev.} \textbf{17} (1975) 443--464.

\bibitem{CsiszarKorner2011}
I.~Csisz\'ar, J.~K\"orner,
\emph{Information Theory: Coding Theorems for Discrete Memoryless Systems},
2nd ed., Cambridge Univ.\ Press (2011), Lemma~2.7.

\bibitem{Audenaert2007}
K.~M.~R.~Audenaert,
``A sharp continuity estimate for the von Neumann entropy,''
\emph{J.\ Phys.\ A} \textbf{40} (2007) 8127--8136;
arXiv:quant-ph/0610146.

\bibitem{Fannes1973}
M.~Fannes,
``A continuity property of the entropy density for spin lattice systems,''
\emph{Commun.\ Math.\ Phys.} \textbf{31} (1973) 291--294.

\bibitem{Seneta2006}
E.~Seneta,
\emph{Non-negative Matrices and Markov Chains},
Springer Series in Statistics, 2nd ed.\ (2006).

\bibitem{LevinPeres2017}
D.~A.~Levin, Y.~Peres,
\emph{Markov Chains and Mixing Times},
2nd ed., American Mathematical Society (2017).

\bibitem{CoverThomas2006}
T.~M.~Cover, J.~A.~Thomas,
\emph{Elements of Information Theory},
2nd ed., Wiley-Interscience (2006).

\bibitem{Ehrlich2026a}
B.~Ehrlich,
``Exponential suppression of transient-basin contributions in trajectory-weighted Markov chain measures,''
preprint (2026).

\bibitem{IpsenMeyer1994}
I.~C.~F.~Ipsen and C.~D.~Meyer,
``Uniform stability of Markov chains,''
\emph{SIAM J.\ Matrix Anal.\ Appl.} \textbf{15} (1994) 1061--1074.

\bibitem{Winter2016}
A.~Winter,
``Tight uniform continuity bounds for quantum entropies: conditional entropy, relative entropy distance, and energy constraints,''
\emph{Commun.\ Math.\ Phys.} \textbf{347} (2016) 291--313.

\end{thebibliography}

\end{document}
